\begin{remark}[3.1.12]
\label{IV.2.3.1.12-fr}
Bajo las hipótesis de (3.1.10), puede ocurrir, en cambio, que $\mathcal F$
carezca de ciclo primo asociado inmerso, pero no así $f_*(\mathcal F)$.
Tomemos, por ejemplo, $Y=\operatorname{Spec}(k[T])$, donde $k$ es un cuerpo
(la ``recta afín''), y sea $X$ la \emph{suma} de $X_1=Y$ y de
$X_2=\operatorname{Spec}(k)$, correspondiendo el morfismo $X_2\to Y$ al
homomorfismo canónico $k[T]\to k[T]/\mathfrak m$, donde $\mathfrak m$ es el
ideal maximal $(T)$. Es claro que el morfismo $f:X\to Y$ es finito; si se toma
$\mathcal F=\mathcal O_X$, $\mathcal F$ carece de ciclo primo asociado
inmerso, pero $f_*(\mathcal F)=\widetilde M$, donde $M$ es el
$k[T]$-módulo suma directa de $k[T]$ y de $k$, de modo que
$\operatorname{Ass}(M)$ está formado por el punto genérico $(0)$ de $Y$ y por
el punto $\mathfrak m$.
\end{remark}

\begin{proposition}[3.1.13]
\label{IV.2.3.1.13-fr}
Sean $X$ un preesquema localmente noetheriano, $U$ un abierto de $X$ e
$i:U\to X$ la inyección canónica. Para todo $\mathcal O_U$-módulo
cuasicoherente $\mathcal F$, se tiene
$\operatorname{Ass}(i_*(\mathcal F))=\operatorname{Ass}(\mathcal F)$.
\end{proposition}

Recordemos que $i_*(\mathcal F)$ es un $\mathcal O_X$-módulo cuasicoherente
(1.2.2 y 1.7.4); como $i_*(\mathcal F)|U=\mathcal F$, se tiene
$(\operatorname{Ass}(i_*(\mathcal F)))\cap U=\operatorname{Ass}(\mathcal F)$,
y queda, por tanto, demostrar que
$\operatorname{Ass}(i_*(\mathcal F))\subset U$. Pero, según (3.1.8), esta
relación significa que, para todo abierto afín $V$ de $X$, toda sección de
$i_*(\mathcal F)$ sobre $V$ que sea nula en $U\cap V$ es nula, condición que
se verifica trivialmente puesto que
$\Gamma(V,i_*(\mathcal F))=\Gamma(U\cap V,\mathcal F)
=\Gamma(U\cap V,i_*(\mathcal F))$.

